% This file should be rename with the speaker's last name. (e.g : smith)

% Write the author names in the order you wish them to appear
% and underline the one who will give the talk.

% Only the speaker's email address is required.

\begin{abstract_online}{The importance of being continuously continuous}{%
        \underline{David Jeffrey}$^{1}$, \underline{David Stoutemyer}$^{2}$}{[djeffrey@uwo.ca; dstout@hawaii.edu]
        }{%
        $^1$ ORCCA, University of Western Ontario, London, Ontario, Canada\newline{}$^2$ Computer Science, University of Hawaii, Hawaii, USA}

We discuss two forms of continuity in the context of integration.

The fundamental theorem of calculus requires that the expression for the integral must be continuous on the interval of interest. Computer Algebra systems, however, do not always co-operate with this requirement. Given an integrand that is continuous on an interval, a computer algebra system may not return an expression that is also continuous on the interval. We show how this can happen, how it can be repaired [1], and speculate on why it has not been.

The other type of continuity refers to continuity with respect to parameters. Consider calculus's most famous integral:

\[ \int x^n\, dx = \frac{x^{n+1}}{n+1} \ .\]

When n=-1, this expression breaks down, but a valid integral still exists, namely log(x). This can be regarded as a discontinuity in the parameter n. There are many similar integrals whose standard expressions contain parametric discontinuities. We show how such integrals can be made parametrically continuous [2], and demonstrate a program that does this.

\textbf{Keywords} \newline{} Integration, Continuity, Kahanian

\textbf{References} \newline{}
% Model for references to books.
[1] \textsc{D.J. Jeffrey}, The importance of being continuous. \emph{Mathematics magazine}, \textbf{67}, 294 - 300, 1994. \newline{}
%  Model for references to articles
[2] \textsc{W. Kahan}, Personal communication. \newline{}
%  Model for references to proceedings.
\end{abstract_online}
